\notechapter{N-20221105}{Function Exercises and Inequality Estimates}

\section{A short note about elementary methods}

A person who understands the structure of a problem sees more than someone who only imitates formulas.  The following exercise sheet develops that principle through periodicity, quadratic minima, logarithmic comparison, parameter ranges, and an inequality with \(abc=1\), with the reasoning made explicit.

\section{Periodicity from a transformed argument}

The first problem has a function \(f\) satisfying a relation of the form
\[
  f(x+4)=f(x),\qquad g(x+4)=g(x)+2,
\]
and asks about the periodicity of a composite or transformed function.  The increment in \(x\) produces the same value of \(f\) and a controlled increment of \(g\).  The general principle is:
\[
  f(x+T)=f(x)\quad\Longrightarrow\quad f\text{ has period }T.
\]
If a formula contains \(x+8\), use the period twice.  For example,
\[
  f(x+8)=f((x+4)+4)=f(x+4)=f(x).
\]

The important point is that a period is not necessarily the least positive period.  A function with period \(4\) automatically has period \(8,12,\ldots\).  Many errors come from confusing ``a period'' with ``the fundamental period.''

\section{A quadratic minimum}

The second problem studies
\[
  y=x^2+x+\frac53.
\]
Complete the square:
\[
  y=\left(x+\frac12\right)^2+\frac{17}{12}.
\]
Hence
\[
  y\ge\frac{17}{12},
\]
with equality at
\[
  x=-\frac12.
\]
This is the cleanest form of a quadratic problem.  Differentiation gives the same result, but completing the square better reveals the geometry: a parabola shifted horizontally and vertically.

\section{A logarithmic comparison}

Another problem compares functions involving logarithms, for example
\[
  f(x)=\log_a(x+1)
\]
and a related expression \(g(x)-f(x)\).  Logarithm rules rewrite the difference as:
\[
  \log_a u-\log_a v=\log_a\frac uv.
\]
Then the sign of the difference depends on whether the base \(a\) is greater than \(1\) or between \(0\) and \(1\).  This is the crucial trap:
\[
  a>1\quad\Rightarrow\quad \log_a x\text{ is increasing},
\]
whereas
\[
  0<a<1\quad\Rightarrow\quad \log_a x\text{ is decreasing}.
\]
So a logarithmic inequality must always track the base.

\section{A parameterized polynomial}

For a parameter \(p\) and a function with a specified zero or extremum, substitute the relevant value and reduce the condition to an equation in \(p\).  The broad template is:
\begin{enumerate}[(i)]
  \item if \(x_0\) is a zero, impose \(f(x_0)=0\);
  \item if \(x_0\) is an extremum of a differentiable function, impose \(f'(x_0)=0\);
  \item if a graph is tangent to a line, impose both value equality and slope equality.

\end{enumerate}

This page is mainly training the habit of translating a verbal condition into equations.

\section[An inequality under abc=1]{An inequality under \(abc=1\)}

The last problem states: for \(a,b,c>0\), \(abc=1\), prove an inequality of the form
\[
  \frac1{a^3(cb+c)}+\frac1{b^3(ca+c)}+\frac1{c^3(a+b)}\ge\frac32.
\]
Use the condition
\[
  abc=1
\]
to replace reciprocal expressions by symmetric sums such as
\[
  \frac1a+\frac1b+\frac1c=ab+bc+ca.
\]
Then Cauchy--Schwarz or AM--GM gives a lower bound.

The useful lesson is that the constraint \(abc=1\) should not sit unused.  It usually allows conversion between variables and reciprocals.  A standard move is also
\[
  a=\frac xy,\qquad b=\frac yz,\qquad c=\frac zx,
\]
which automatically enforces \(abc=1\).  Whether one uses this substitution or direct reciprocal rewriting depends on the shape of the expression.

\section{The conceptual turn}

The exercises look miscellaneous, but the same strategy recurs: normalize the expression before attacking it.  Periodicity reduces arguments modulo a period; completing the square normalizes a quadratic; logarithm laws turn products into sums; parameter conditions become equations; \(abc=1\) turns reciprocal expressions into symmetric ones.  Elementary problem solving is often the art of finding the right normalization.

\section{Periodicity from transformed arguments}

A relation such as \(f(x+a)=f(x)\) is invariance under translation by \(a\).  More complicated relations, for example involving \(f(c-x)\), combine translation and reflection.  The group generated by these transformations controls the function's repeated values.

Thus many functional-equation exercises are really symmetry problems on the domain.

\section{Quadratic and logarithmic estimates}

A quadratic minimum is controlled by completing the square:
\[
  ax^2+bx+c=a\left(x+\frac b{2a}\right)^2+c-\frac{b^2}{4a}.
\]
A logarithmic comparison often uses monotonicity or concavity of \(\log\).  In both cases, the point is to expose the structural feature: convexity for quadratics, monotone concavity for logarithms.

\section{Functional equations and orbit reduction}

If a functional equation lets one transform \(x\) into simpler inputs, then the domain splits into orbits under the generated transformations.  A fundamental domain is a set of representatives for these orbits.

Solving the equation means assigning values on the fundamental domain and transporting them consistently along the orbit.

\section{Functions through symmetry and normal form}

Function exercises are unified by structure: periodicity is group invariance, quadratics are convex normal forms, logarithmic inequalities come from monotonicity and concavity, and functional equations reduce values along orbits.
